Zu zeigen ist, dass die kinetische Energie vor und nach dem Stoss gleich ist, also
\begin{align}
 m_1\,{v_1'}^2 + m_2\,{v_2'}^2 &= m_1\,v_1^2 + m_2\,v_2^2
\end{align}
(der Faktor $\frac{1}{2}$ fällt auf beiden Seiten weg). Mit $M = m_1+m_2$ multiplizieren wir die linke Seite mit $M^2$ und setzen die Formeln ein:
\begin{align}
 M^2\left(m_1\,{v_1'}^2 + m_2\,{v_2'}^2\right) &= m_1\left[\left(m_1-m_2\right)v_1+2m_2\,v_2\right]^2 + m_2\left[\left(m_2-m_1\right)v_2+2m_1\,v_1\right]^2
\end{align}
Wir sortieren die ausmultiplizierten Terme nach $v_1^2$, $v_1 v_2$ und $v_2^2$:
\begin{align}
 v_1^2:&\quad m_1\left(m_1-m_2\right)^2 + 4m_1^2 m_2 = m_1\left[\left(m_1-m_2\right)^2+4m_1 m_2\right] = m_1\left(m_1+m_2\right)^2 = m_1 M^2 \\
 v_1 v_2:&\quad 4m_1 m_2\left(m_1-m_2\right) + 4m_1 m_2\left(m_2-m_1\right) = 0 \\
 v_2^2:&\quad 4m_1 m_2^2 + m_2\left(m_2-m_1\right)^2 = m_2\left[\left(m_1-m_2\right)^2+4m_1 m_2\right] = m_2\left(m_1+m_2\right)^2 = m_2 M^2
\end{align}
Also ist
\begin{align}
 M^2\left(m_1\,{v_1'}^2 + m_2\,{v_2'}^2\right) &= M^2\left(m_1\,v_1^2 + m_2\,v_2^2\right)
\end{align}
und nach Division durch $M^2$ und Multiplikation mit $\frac{1}{2}$
\begin{align}
 \result{\frac{1}{2}m_1\,{v_1'}^2 + \frac{1}{2}m_2\,{v_2'}^2 = \frac{1}{2}m_1\,v_1^2 + \frac{1}{2}m_2\,v_2^2} \quad\square
\end{align}
